\providecommand{\click}[2]{\href{panel:#1}{#2}}
\providecommand{\PanelAnim}[1]{\href{anim:#1}{ }}
\providecommand{\PanelTitle}[1]{\section*{#1}}


\PanelTitle{Pendulum phase portrait and stability}

We consider the nonlinear pendulum

\[
\dot{x}_1 = x_2, \qquad
\dot{x}_2 = -a\sin x_1 - b x_2 .
\]

The variable \(x_1\) is the angle of the pendulum and \(x_2\) is its
angular velocity. Each point in the plane \((x_1,x_2)\) represents a
possible state of the system.

The phase portrait shows how the state evolves. A curve represents the
trajectory generated by some initial condition.

Equilibrium points occur when

\[
x_2 = 0, \qquad \sin x_1 = 0 .
\]

Hence equilibria appear at

\[
(x_1,x_2) = (k\pi,0), \qquad k\in\mathbb{Z}.
\]

In the window displayed below we see two of them:

\[
(0,0) \quad \text{and} \quad (\pi,0).
\]

The point \(x_1=0\) corresponds to the pendulum hanging downward.
The point \(x_1=\pi\) corresponds to the pendulum balanced upside down.

The two orange circles represent small neighborhoods around the two
equilibria. In the language of Lyapunov stability, these circles play
the role of an \(\varepsilon\)-ball around the equilibrium.

\PanelAnim{pendulum}

\textbf{Understanding Lyapunov stability}

Recall the definition of stability:

\[
\|x(0)\| < \delta
\quad\Rightarrow\quad
\|x(t)\| < \varepsilon
\quad \forall t \ge 0 .
\]

This means that trajectories starting sufficiently close to the
equilibrium must remain inside the \(\varepsilon\)-ball forever.

The circles in the figure help you test this idea visually.

\textbf{Experiment 1: stability without friction}

Set

\[
b = 0 .
\]

Click inside the circle centered at \((0,0)\).

You should observe closed curves surrounding the equilibrium.
Trajectories never leave the circle if you start sufficiently close to
the origin.

However they do not converge to the equilibrium. Instead they remain on
periodic orbits.

This illustrates an important idea: the equilibrium is stable but not
asymptotically stable.

\textbf{Experiment 2: asymptotic stability with damping}

Increase \(b\).

Now click again inside the left circle.

Trajectories spiral toward the origin. Once they enter the circle they
remain inside and eventually approach the equilibrium.

This behavior corresponds to asymptotic stability: trajectories not
only stay near the equilibrium, but converge to it as \(t\to\infty\).

Physically, friction dissipates energy and the pendulum eventually
comes to rest.

\textbf{Experiment 3: the unstable inverted pendulum}

Now click inside the circle centered at

\[
(\pi,0).
\]

Even when the initial condition is extremely close to this point, the
trajectory quickly leaves the circle.

This violates the Lyapunov stability condition: there is always a
trajectory that escapes the \(\varepsilon\)-ball.

Therefore the inverted pendulum equilibrium is unstable.

\textbf{Experiment 4: saddle behavior}

Try several initial conditions near \((\pi,0)\).

You will notice that trajectories approach the equilibrium only along
very special directions and escape along others.

This is the typical behavior of a saddle point.

\textbf{Key geometric idea}

Lyapunov stability is a geometric property of trajectories.

Stable equilibrium:
trajectories starting inside a small neighborhood remain inside.

Unstable equilibrium:
some trajectories leave the neighborhood no matter how small it is.

The phase portrait allows us to observe these properties directly,
without solving the differential equation explicitly.
